\section{September 30, 2026}
% Based on pages 1--3 of 6-7-week.pdf.

We continue the discussion of time reversibility from September 28.
Symmetry constrains the error expansion and makes extrapolation
especially effective. We also introduce variable doubling and the
perturbation viewpoint behind modified equations.
Throughout, \(\Phi\) and \(\Psi\) denote the exact and numerical
flows, and \(\tau\) denotes a constant step size.

\subsection{Reversible integrators and explicit RK methods}

Recall that a one-step method is symmetric if
\[
 \Psi^{t,t+\tau}\circ\Psi^{t+\tau,t}=I,
\]
where the first map uses the same method with step \(-\tau\).
Gauss collocation has this property because its nodes are invariant
under reflection about the midpoint of the step.

\begin{proposition}
 A consistent explicit RK method with finitely many stages cannot
 be symmetric for general ODEs.
\end{proposition}

\begin{proof}
 Apply the method to \(x'=x\), with initial value one. Its stability
 function \(R(z)\) is a polynomial, and consistency gives
 \(R(0)=R'(0)=1\). A forward step followed by a backward step
 multiplies the initial value by \(R(-\tau)R(\tau)\).
 Symmetry would imply
 \[
  R(-z)R(z)=1
 \]
 for all sufficiently small real \(z\), hence as a polynomial
 identity. If \(R\) has degree \(m\geq1\) and leading coefficient
 \(a_m\), the product has degree \(2m\) and leading coefficient
 \((-1)^m a_m^2\neq0\). It cannot equal the constant one.
\end{proof}

This obstruction concerns standard explicit RK methods on a general
vector field. A method exploiting a partition of the variables can
still be both explicit and reversible.

\subsection{Variable doubling}

For \(x'=f(t,x)\), introduce a second copy of the unknown and consider
\begin{equation}\label{eq:doubled-system-sep-30}
 x'=f(t,y),\qquad y'=f(t,x),\qquad x(t_0)=y(t_0)=x_0.
\end{equation}
The diagonal \(x=y\) is invariant: the pair consisting of two copies
of the original solution solves this system, and uniqueness identifies
it with the solution of \eqref{eq:doubled-system-sep-30}.

The following staggered update is explicit:
\begin{align}
 y_{n+1/2}&=y_n+\frac\tau2 f(t_n,x_n),\notag\\
 x_{n+1}&=x_n+\tau f(t_n+\tau/2,y_{n+1/2}),
       \label{eq:doubled-step-sep-30}\\
 y_{n+1}&=y_{n+1/2}+\frac\tau2 f(t_n+\tau,x_{n+1}).\notag
\end{align}
Equivalently,
\begin{align*}
 x_{n+1}&=x_n+\tau f\left(t_n+\frac\tau2,
                         y_n+\frac\tau2 f(t_n,x_n)\right),\\
 y_{n+1}&=y_n+\frac\tau2
           \bigl(f(t_n,x_n)+f(t_n+\tau,x_{n+1})\bigr).
\end{align*}
Here \(y_{n+1/2}\) is an intermediate variable, not an additional
solution value required from a previous step.

To reverse the step, start from \((x_{n+1},y_{n+1})\) at
\(t_n+\tau\) and replace \(\tau\) by \(-\tau\). The three
operations successively recover
\[
 y_{n+1}-\frac\tau2 f(t_n+\tau,x_{n+1})=y_{n+1/2},
 \qquad x_n,\qquad y_n.
\]
Thus the map on the pair \((x,y)\) is reversible.

Taylor expansion also shows that it has order two for the doubled
system. For example, its first component satisfies
\[
 x_{n+1}=x_n+\tau f(t_n,y_n)
 +\frac{\tau^2}{2}
 \bigl(f_t(t_n,y_n)+f_x(t_n,y_n)f(t_n,x_n)\bigr)
 +O(\tau^3),
\]
which agrees with the exact expansion of \(x'=f(t,y)\).
The corresponding statement holds for the second component.
Numerically, \(x_n\) and \(y_n\) need not be equal. Both copies must
be retained between steps for the reversibility argument to apply;
resetting one copy to the other defines a different method.

\subsection{Even powers in the global error expansion}

Write \(x_\tau(t_n)=x_n\) for the numerical solution computed with
constant step \(\tau\) and exact initial data. At a fixed endpoint,
\(x_\tau(T)=Y_\tau(T)\) in the notation of September 9.
Comparisons between step sizes use common grid points on a fixed
finite interval.

\begin{theorem}[Error expansion for symmetric methods]
 Suppose a smooth, consistent symmetric one-step method has finite
 order and admits the usual asymptotic error expansion. Its order
 is even, say \(2r\), and, with sufficient smoothness,
 \begin{equation}\label{eq:even-error-sep-30}
  x_\tau(t)=x(t)+C_{2r}(t)\tau^{2r}
       +C_{2r+2}(t)\tau^{2r+2}+\cdots.
 \end{equation}
 The expansion is understood to any finite order for which the
 required derivatives exist; no convergence of the infinite series
 is asserted.
\end{theorem}

One explanation uses the formal modified equation developed next
lecture. In the autonomous setting, write the numerical map formally
as the time-\(\tau\) flow of a step-dependent vector field
\(g_\tau\). Symmetry says that its inverse is the same map with
step \(-\tau\). Taking the formal logarithm gives
\[
 -\tau g_{-\tau}=-\tau g_\tau,
 \qquad g_{-\tau}=g_\tau.
\]
Thus only even powers of \(\tau\) occur in \(g_\tau\), and hence
in the fixed-time solution expansion. A nonautonomous equation can
be treated by adjoining \(t'=1\). The statement requires the fixed
step and initialization assumptions above; arbitrary interpolation
between grid points or arbitrary step changes do not automatically
preserve this expansion.

\subsection{Richardson extrapolation in the squared step size}

We specialize \eqref{eq:richardson-sep-09} to a second-order
symmetric method:
\begin{align*}
 x_\tau(t)&=x(t)+C_2(t)\tau^2+C_4(t)\tau^4+O(\tau^6),\\
 x_{\tau/2}(t)&=x(t)+\frac14C_2(t)\tau^2
                          +\frac1{16}C_4(t)\tau^4+O(\tau^6).
\end{align*}
Eliminating the leading error gives
\begin{equation}\label{eq:richardson-four-sep-30}
 \boxed{\frac{4x_{\tau/2}(t)-x_\tau(t)}3
        =x(t)-\frac14C_4(t)\tau^4+O(\tau^6).}
\end{equation}
One extrapolation therefore raises the order from two to four.
For an expansion beginning at order \(2r\), the corresponding
combination is
\[
 \frac{2^{2r}x_{\tau/2}(t)-x_\tau(t)}{2^{2r}-1}
       =x(t)+O(\tau^{2r+2}).
\]

Using three step sizes, fit a quadratic polynomial in \(s=\tau^2\):
\[
 P(s)=d_0+d_2s+d_4s^2,
\]
with values \(x_\tau(t),x_{\tau/2}(t),x_{\tau/4}(t)\) at
\(s=\tau^2,\tau^2/4,\tau^2/16\). Componentwise, the system is
\[
 \begin{bmatrix}
  1&\tau^2&\tau^4\\
  1&\tau^2/4&\tau^4/16\\
  1&\tau^2/16&\tau^4/256
 \end{bmatrix}
 \begin{bmatrix}d_0\\d_2\\d_4\end{bmatrix}
 =\begin{bmatrix}x_\tau(t)\\x_{\tau/2}(t)\\x_{\tau/4}(t)\end{bmatrix}.
\]
Evaluating the interpolant at zero gives
\begin{equation}\label{eq:richardson-six-sep-30}
 \boxed{d_0=\frac1{45}x_\tau(t)-\frac{20}{45}x_{\tau/2}(t)
                     +\frac{64}{45}x_{\tau/4}(t)
             =x(t)+O(\tau^6).}
\end{equation}
The weights sum to one and annihilate the \(\tau^2\) and
\(\tau^4\) terms. This improves accuracy at the comparison time;
a linear combination of numerical solutions need not retain the
geometric properties of the underlying symmetric method.

\subsection{Perturbation expansions}

As a model for coefficient matching, consider weak damping:
\[
 q''=-q-2\varepsilon q',\qquad
 q(t)=q_0(t)+\varepsilon q_1(t)+\varepsilon^2q_2(t)+\cdots.
\]
Equating successive powers of \(\varepsilon\) gives
\[
 q_0''=-q_0,\qquad
 q_1''=-q_1-2q_0',\qquad
 q_m''=-q_m-2q_{m-1}'\quad(m\geq1).
\]
For initial data independent of \(\varepsilon\), the higher
corrections have zero initial position and velocity. For example,
if \(q(0)=1\) and \(q'(0)=0\), then
\[
 q_0(t)=\cos t,\qquad q_1(t)=\sin t-t\cos t.
\]
The secular term \(t\cos t\) also reminds us that a fixed-time
asymptotic expansion need not be uniform for arbitrarily long times.

For a numerical method, the small parameter is the step size.
Forward error analysis expands its solution around the exact one;
backward error analysis instead seeks a nearby differential equation
whose flow agrees with the numerical method to a prescribed order.
This leads to modified equations and, for gradient flows, an
interpretation in terms of a modified potential.
